\section{The Second Fundamental Form}

Let $(\widetilde M,\widetilde g)$ be an $m$-dimensional Riemannian or
pseudo-Riemannian manifold without boundary, and let $(M,g)$ be an embedded
$n$-dimensional Riemannian submanifold without boundary, with
$g=\iota_M^*\widetilde g$. All statements are local and therefore also apply
to isometric immersions after restriction to embedded neighborhoods. We
restrict to positive-definite induced metrics; mixed-signature submanifolds
require additional qualifications (see \cite{dC92}).

Ambient covariant derivatives and curvatures carry tildes; untilded symbols
refer to the induced geometry of $M$. The orthogonal splitting
$T\widetilde M|_M=TM\oplus NM$ has smooth projections
\[
\pi^\top:T\widetilde M|_M\to TM,
\qquad
\pi^\perp:T\widetilde M|_M\to NM.
\]
For a section $A$ of $T\widetilde M|_M$, write
$A^\top=\pi^\top A$ and $A^\perp=\pi^\perp A$. If
$X,Y\in\mathfrak X(M)$ are extended to ambient vector fields, then
\begin{equation}
\widetilde\nabla_XY
=(\widetilde\nabla_XY)^\top+(\widetilde\nabla_XY)^\perp.
\tag{8.1}
\end{equation}
The \emph{second fundamental form} is the normal-bundle-valued map
\[
\mathrm{II}(X,Y)=(\widetilde\nabla_XY)^\perp.
\]

\begin{proposition}[Properties of the Second Fundamental Form]
\label{prop:second-fundamental-form-properties}
\dcref{ch8:8.1}

Let $(M,g)$ be an embedded Riemannian submanifold of a Riemannian or
pseudo-Riemannian manifold $(\widetilde M,\widetilde g)$. For
$X,Y\in\mathfrak X(M)$:
\begin{enumerate}[label=(\alph*)]
\item $\mathrm{II}(X,Y)$ is independent of the chosen ambient extensions.
\item $\mathrm{II}$ is $C^\infty(M)$-bilinear.
\item $\mathrm{II}(X,Y)=\mathrm{II}(Y,X)$.
\item $\mathrm{II}(X,Y)|_p$ depends only on $X_p$ and $Y_p$.
\end{enumerate}
\end{proposition}

\begin{proof}
\sketch For fixed ambient extensions, symmetry of $\widetilde\nabla$ gives
\[
\mathrm{II}(X,Y)-\mathrm{II}(Y,X)
=(\widetilde\nabla_XY-\widetilde\nabla_YX)^\perp=[X,Y]^\perp=0,
\]
because the bracket of fields tangent along $M$ is tangent to $M$. At
$p\in M$, $\widetilde\nabla_XY|_p$ depends only on $X_p$, proving independence
of the extension of $X$ and $C^\infty(M)$-linearity in $X$. Symmetry gives the
same conclusions for $Y$, and hence pointwise dependence on both arguments.
\end{proof}

Thus $\mathrm{II}(v,w)$ is well defined for $v,w\in T_pM$ by extending them
to tangent vector fields.

\begin{theorem}[The Gauss Formula]
\label{thm:gauss-formula}
\uses{prop:second-fundamental-form-properties, ex:tangential-connection-submanifold, prop:tangential-connection-compatible-metric, prop:tangential-connection-symmetric}
\dcref{ch8:8.2}


Under the preceding hypotheses, for arbitrary ambient extensions of
$X,Y\in\mathfrak X(M)$,
\[
\widetilde\nabla_XY=\nabla_XY+\mathrm{II}(X,Y)
\qquad\text{along }M.
\]
\end{theorem}

\begin{proof}
\sketch Define
$\nabla_X^\top Y=(\widetilde\nabla_XY)^\top$. The extension-independence
argument for the second fundamental form shows that $\nabla^\top$ is a
connection on $M$. Ambient metric compatibility and orthogonal projection
make it compatible with $g$, while ambient symmetry and tangency of brackets
make it symmetric. Uniqueness of the Levi-Civita connection gives
$\nabla^\top=\nabla$. Equation (8.1) then proves the formula.
\end{proof}

\begin{corollary}[The Gauss Formula Along a Curve]
\label{cor:gauss-formula-along-curve}
\uses{thm:gauss-formula}
\dcref{ch8:8.3}


Let $\gamma:I\to M$ be smooth and let $X$ be a smooth vector field along
$\gamma$ tangent to $M$. Its ambient and intrinsic covariant derivatives
satisfy
\begin{equation}
\widetilde D_tX=D_tX+\mathrm{II}(\gamma',X).
\tag{8.2}
\end{equation}
\end{corollary}

\begin{proof}
\sketch Near $\gamma(t_0)$ choose an adapted orthonormal frame
$(E_1,\ldots,E_m)$ and write $X=X^iE_i$ for $1\leq i\leq n$. The product rule
and Gauss formula give
\[
\widetilde D_tX
=\sum_{i=1}^n\left(\dot X^iE_i+X^i\nabla_{\gamma'}E_i
+X^i\mathrm{II}(\gamma',E_i)\right)
=D_tX+\mathrm{II}(\gamma',X).
\]
\end{proof}

For $N\in\Gamma(NM)$, define the symmetric scalar-valued form
\begin{equation}
\mathrm{II}_N(X,Y)=\langle N,\mathrm{II}(X,Y)\rangle.
\tag{8.3}
\end{equation}
Its associated self-adjoint bundle map $W_N:TM\to TM$, the
\emph{Weingarten map in the direction $N$}, is characterized by
\begin{equation}
\langle W_N(X),Y\rangle
=\mathrm{II}_N(X,Y)
=\langle N,\mathrm{II}(X,Y)\rangle.
\tag{8.4}
\end{equation}

\begin{proposition}[The Weingarten Equation]
\label{prop:weingarten-equation}
\dcref{ch8:8.4}

For every $X\in\mathfrak X(M)$ and $N\in\Gamma(NM)$,
\begin{equation}
(\widetilde\nabla_XN)^\top=-W_N(X),
\tag{8.5}
\end{equation}
where $N$ is extended arbitrarily to a neighborhood of $M$ in
$\widetilde M$.
\end{proposition}

\begin{proof}
\sketch For arbitrary $Y\in\mathfrak X(M)$, metric compatibility, normality
of $N$, and the Gauss formula yield
\begin{align*}
0=X\langle N,Y\rangle
&=\langle\widetilde\nabla_XN,Y\rangle
 +\langle N,\widetilde\nabla_XY\rangle\\
&=\langle\widetilde\nabla_XN,Y\rangle
 +\langle N,\mathrm{II}(X,Y)\rangle\\
&=\langle(\widetilde\nabla_XN)^\top+W_N(X),Y\rangle.
\end{align*}
Since $Y$ is arbitrary, (8.5) follows.
\end{proof}

\begin{theorem}[The Gauss Equation]
\label{thm:gauss-equation}
\uses{prop:covariant-derivative-depends-on-curve}
\dcref{ch8:8.5}


For all $W,X,Y,Z\in\mathfrak X(M)$,
\[
\widetilde{Rm}(W,X,Y,Z)
=Rm(W,X,Y,Z)
-\langle\mathrm{II}(W,Z),\mathrm{II}(X,Y)\rangle
+\langle\mathrm{II}(W,Y),\mathrm{II}(X,Z)\rangle.
\]
\end{theorem}

\begin{proof}
\sketch Extend $W,X,Y,Z$ ambiently. Apply the Gauss formula inside the
definition of ambient curvature:
\begin{align*}
\widetilde{Rm}(W,X,Y,Z)
&=\big\langle
\widetilde\nabla_W(\nabla_XY+\mathrm{II}(X,Y))
-\widetilde\nabla_X(\nabla_WY+\mathrm{II}(W,Y))\\
&\qquad-\widetilde\nabla_{[W,X]}Y,Z
\big\rangle.
\end{align*}
The Weingarten equation converts the tangential components of the two normal
derivatives into
$-\langle\mathrm{II}(X,Y),\mathrm{II}(W,Z)\rangle$ and
$+\langle\mathrm{II}(W,Y),\mathrm{II}(X,Z)\rangle$. Applying the Gauss
formula to the remaining tangential derivatives leaves
$\langle R(W,X)Y,Z\rangle$, proving the equation.
\end{proof}

The \emph{normal connection} is
\[
\nabla_X^\perp N=(\widetilde\nabla_XN)^\perp
\qquad
(X\in\mathfrak X(M),\ N\in\Gamma(NM)).
\]

\begin{proposition}
\label{prop:normal-connection-well-defined}
\dcref{ch8:8.6}

The operator $\nabla^\perp$ is a well-defined connection on $NM$ compatible
with the restricted ambient metric:
\[
X\langle N_1,N_2\rangle
=\langle\nabla_X^\perp N_1,N_2\rangle
 +\langle N_1,\nabla_X^\perp N_2\rangle
\]
for $N_1,N_2\in\Gamma(NM)$ and $X\in\mathfrak X(M)$.
\end{proposition}

Let $F\to M$ be the bundle with fiber
\[
F_p=\operatorname{Bil}(T_pM\times T_pM,N_pM).
\]
For $B\in\Gamma(F)$, define
\[
(\nabla_X^FB)(Y,Z)
=\nabla_X^\perp(B(Y,Z))
-B(\nabla_XY,Z)-B(Y,\nabla_XZ).
\]

\begin{theorem}[The Codazzi Equation]
\label{thm:codazzi-equation}
\uses{thm:gauss-equation}
\dcref{ch8:8.9}


For all $W,X,Y\in\mathfrak X(M)$,
\begin{equation}
(\widetilde R(W,X)Y)^\perp
=(\nabla_W^F\mathrm{II})(X,Y)-(\nabla_X^F\mathrm{II})(W,Y).
\tag{8.6}
\end{equation}
\end{theorem}

\begin{proof}
\sketch It suffices to pair both sides with an arbitrary normal field $N$:
\begin{equation}
\langle\widetilde R(W,X)Y,N\rangle
=\langle(\nabla_W^F\mathrm{II})(X,Y),N\rangle
-\langle(\nabla_X^F\mathrm{II})(W,Y),N\rangle.
\tag{8.7}
\end{equation}
After applying the Gauss formula to the ambient curvature definition, retain
only normal components. The first ambient derivative contributes
\[
\mathrm{II}(W,\nabla_XY)
+(\nabla_W^F\mathrm{II})(X,Y)
+\mathrm{II}(\nabla_WX,Y)
+\mathrm{II}(X,\nabla_WY),
\]
and the second contributes the analogous expression with $W,X$ interchanged,
while the bracket term contributes $-\mathrm{II}([W,X],Y)$. The terms
containing $\nabla_XY$ and $\nabla_WY$ cancel, and
\[
\mathrm{II}(\nabla_WX-\nabla_XW-[W,X],Y)=0
\]
by symmetry of $\nabla$. Pairing the remaining terms with $N$ gives (8.7).
\end{proof}

\section{Curvature of a Curve}

For a smooth unit-speed curve $\gamma:I\to M$ in a Riemannian or
pseudo-Riemannian manifold, its \emph{geodesic curvature} is
\[
\kappa(t)=|D_t\gamma'(t)|.
\]
For a regular curve in a Riemannian manifold, define curvature using any
unit-speed reparametrization. In the pseudo-Riemannian case, this definition
is restricted to curves satisfying $|\gamma'(t)|\neq0$ everywhere. A
unit-speed curve has zero geodesic curvature exactly when it is a geodesic.

If $M$ is a Riemannian submanifold of
$(\widetilde M,\widetilde g)$, a regular curve in $M$ has intrinsic curvature
$\kappa$ computed in $M$ and extrinsic curvature $\widetilde\kappa$ computed
in $\widetilde M$.

\begin{proposition}[Geometric Interpretation of II]
\label{prop:geometric-interpretation-ii}
\uses{cor:gauss-formula-along-curve}
\dcref{ch8:8.10}


Let $p\in M$ and $v\in T_pM$.
\begin{enumerate}
\item[(a)] $\mathrm{II}(v,v)$ is the ambient acceleration at $p$ of the
$g$-geodesic $\gamma_v$.
\item[(b)] If $v$ is a unit vector, then $|\mathrm{II}(v,v)|$ is the ambient
curvature of $\gamma_v$ at $p$.
\end{enumerate}
\end{proposition}

\begin{proof}
\sketch For any regular $\gamma:(-\varepsilon,\varepsilon)\to M$ with
$\gamma(0)=p$ and $\gamma'(0)=v$, (8.2) gives
\[
\widetilde D_t\gamma'
=D_t\gamma'+\mathrm{II}(\gamma',\gamma').
\]
For $\gamma=\gamma_v$, the intrinsic derivative vanishes. If $v$ is unit,
the ambient acceleration norm is the ambient curvature.
\end{proof}

\begin{lemma}
\label{lem:symmetric-bilinear-form-unit-vectors}
\dcref{ch8:8.11}

Let $V$ be an inner product space, $W$ a vector space, and
$B,B':V\times V\to W$ symmetric bilinear maps. If
$B(v,v)=B'(v,v)$ for every unit vector $v$, then $B=B'$.
\end{lemma}

\begin{proof}
\sketch Homogeneity extends the diagonal equality to every $v\in V$.
Polarization then gives
\[
B(v,w)=\frac14\bigl(B(v+w,v+w)-B(v-w,v-w)\bigr),
\]
and the same identity for $B'$ proves equality.
\end{proof}

A Riemannian submanifold $(M,g)$ of $(\widetilde M,\widetilde g)$ is
\emph{totally geodesic} if every ambient geodesic tangent to $M$ at a time
$t_0$ remains in $M$ on some interval about $t_0$.

\begin{proposition}
\label{prop:totally-geodesic-equivalence}
\uses{thm:gauss-formula, lem:symmetric-bilinear-form-unit-vectors}
\dcref{ch8:8.12}


The following are equivalent:
\begin{enumerate}
\item[(a)] $M$ is totally geodesic in $\widetilde M$.
\item[(b)] Every $g$-geodesic in $M$ is an ambient
$\widetilde g$-geodesic.
\item[(c)] $\mathrm{II}=0$ identically.
\end{enumerate}
\end{proposition}

\begin{proof}
\sketch Assume (a), let $\gamma$ be a $g$-geodesic, and at each $t_0$ let
$\widetilde\gamma$ be the ambient geodesic with the same initial data.
Locally $\widetilde\gamma$ lies in $M$, and (8.2) decomposes
$0=\widetilde D_t\widetilde\gamma'$ into tangent and normal parts. Thus it is
also a $g$-geodesic, so uniqueness gives
$\widetilde\gamma=\gamma$. Hence (b).

Under (b), for the $g$-geodesic with initial vector $v\in T_pM$, both
intrinsic and ambient accelerations vanish. Formula (8.2) gives
$\mathrm{II}(v,v)=0$, and
Lemma~\ref{lem:symmetric-bilinear-form-unit-vectors} gives (c).

If (c) holds, every $g$-geodesic is ambient by (8.2). Given an ambient
geodesic initially tangent to $M$, compare it with the $g$-geodesic having
the same initial data. Ambient geodesic uniqueness makes them agree near the
initial time, proving (a).
\end{proof}
\section{Hypersurfaces}

Throughout these sections, $(M,g)$ is an embedded Riemannian hypersurface of
dimension $n$ in a Riemannian manifold $(\widetilde M,\widetilde g)$ of
dimension $n+1$.

\section{The Scalar Second Fundamental Form and the Shape Operator}

Locally, choose a smooth unit normal field $N$ along $M$. Such a choice always
exists near each point; if both $M$ and $\widetilde M$ are orientable, their
orientations determine a global choice. The \emph{scalar second fundamental
form} relative to $N$ is the symmetric covariant tensor
\begin{equation}
h(X,Y)=\langle N,\operatorname{II}(X,Y)\rangle. \tag{8.8}
\end{equation}
The Gauss formula and orthogonality give
\begin{equation}
h(X,Y)=\langle N,\widetilde\nabla_XY\rangle, \tag{8.9}
\end{equation}
and, because $N$ spans the normal bundle,
\begin{equation}
\operatorname{II}(X,Y)=h(X,Y)N. \tag{8.10}
\end{equation}
Replacing $N$ by $-N$ changes $h$ to $-h$.

The \emph{shape operator} $s=W_N$ is obtained from $h$ by raising one index:
\[
\langle sX,Y\rangle=h(X,Y).
\]
It is self-adjoint, and replacing $N$ by $-N$ changes $s$ to $-s$.

For symmetric covariant $2$-tensors $h,k$, use the Kulkarni--Nomizu product
\begin{align*}
(h\owedge k)(w,x,y,z)
&=h(w,z)k(x,y)+h(x,y)k(w,z)\\
&\quad-h(w,y)k(x,z)-h(x,z)k(w,y),
\end{align*}
and for a symmetric covariant $2$-tensor $T$, use the exterior covariant
derivative
\[
(DT)(x,y,z)=-(\nabla T)(x,y,z)+(\nabla T)(x,z,y).
\]

\begin{theorem}[Fundamental Equations for a Hypersurface]
\label{thm:fundamental-equations-hypersurface}
\dcref{ch8:8.13}
Let $(M,g)$ be a Riemannian hypersurface in $(\overline M,\overline g)$, and
let $N$ be a smooth unit normal field along $M$.
\begin{enumerate}
\item[(a)] If $X,Y\in\mathfrak X(M)$ are extended to an open subset of
$\overline M$, then
\[
\overline\nabla_XY=\nabla_XY+h(X,Y)N.
\]
\item[(b)] If $\gamma\colon I\to M$ is smooth and $X$ is a smooth tangent
vector field along $\gamma$, then
\[
\overline D_tX=D_tX+h(\gamma',X)N.
\]
\item[(c)] For every $X\in\mathfrak X(M)$,
\begin{equation}
\overline\nabla_XN=-sX. \tag{8.11}
\end{equation}
\item[(d)] For all $W,X,Y,Z\in\mathfrak X(M)$,
\[
\overline{\operatorname{Rm}}(W,X,Y,Z)
=\operatorname{Rm}(W,X,Y,Z)-\frac12(h\owedge h)(W,X,Y,Z).
\]
\item[(e)] For all $W,X,Y\in\mathfrak X(M)$,
\begin{equation}
\overline{\operatorname{Rm}}(W,X,Y,N)=(Dh)(Y,W,X). \tag{8.12}
\end{equation}
\end{enumerate}
\end{theorem}

\begin{proof}
\sketch Substitution of (8.10) into the general Gauss formulas and Gauss
equation gives (a), (b), and (d). The general Weingarten equation gives
$(\overline\nabla_XN)^\top=-sX$, while
\[
\langle\overline\nabla_XN,N\rangle
=\frac12X(|N|^2)=0;
\]
hence $\overline\nabla_XN$ is tangent and (c) follows.

For (e), the unit-length condition gives
$\nabla_X^\perp N=0$, since $N$ spans the normal bundle. Therefore
\begin{align*}
(\nabla_W^\perp\operatorname{II})(X,Y)
&=\nabla_W^\perp(h(X,Y)N)
  -\operatorname{II}(\nabla_WX,Y)
  -\operatorname{II}(X,\nabla_WY)\\
&=(\nabla_Wh)(X,Y)N.
\end{align*}
The general Codazzi equation then yields
\begin{align*}
\overline{\operatorname{Rm}}(W,X,Y,N)
&=(\nabla_Wh)(X,Y)-(\nabla_Xh)(W,Y)\\
&=(\nabla h)(Y,X,W)-(\nabla h)(Y,W,X)\\
&=(Dh)(Y,W,X),
\end{align*}
using symmetry of $\nabla h$ in its first two tensor arguments.
\end{proof}

\section{Principal Curvatures}

\begin{lemma}
\label{lem:self-adjoint-eigenvector-extremum}
\dcref{ch8:8.14}
Let $V$ be a finite-dimensional inner product space, let
$s\colon V\to V$ be self-adjoint, and let $C$ be the unit sphere in $V$.
The function $v\mapsto\langle sv,v\rangle$ attains its maximum at some
$v_0\in C$, and every maximizer is an eigenvector of $s$ with eigenvalue
$\lambda_0=\langle sv_0,v_0\rangle$.
\end{lemma}

\begin{proposition}[Finite-Dimensional Spectral Theorem]
\label{prop:finite-dimensional-spectral-theorem}
\uses{lem:self-adjoint-eigenvector-extremum}
\dcref{ch8:8.16}

Every self-adjoint endomorphism of a finite-dimensional inner product space
has an orthonormal basis of eigenvectors, and all its eigenvalues are real.
\end{proposition}

\begin{proof}
\sketch Induct on $n=\dim V$. The result is immediate for $n=1$. In dimension
$n+1$, the preceding lemma supplies a unit eigenvector $b_0$ with real
eigenvalue. For $B=\operatorname{span}\{b_0\}$, self-adjointness implies
$s(B^\perp)\subseteq B^\perp$. Apply the induction hypothesis to
$s|_{B^\perp}$ and adjoin $b_0$ to the resulting orthonormal eigenbasis.
\end{proof}

At $p\in M$, let $\kappa_1,\ldots,\kappa_n$ be the eigenvalues of the shape
operator. In an orthonormal eigenbasis $(b_i)$,
\[
sb_i=\kappa_i b_i,
\qquad
h(v,w)=\sum_{i=1}^n\kappa_i v^iw^i.
\]
The $\kappa_i$ are the \emph{principal curvatures}, and the corresponding
eigenspaces are the \emph{principal directions}. Reversing $N$ negates every
principal curvature but leaves the principal directions unchanged.

The \emph{Gaussian curvature} and \emph{mean curvature} are
\[
K=\det s=\prod_{i=1}^n\kappa_i,
\qquad
H=\frac1n\operatorname{tr}s
=\frac1n\operatorname{tr}_g h
=\frac1n\sum_{i=1}^n\kappa_i.
\]
Under $N\mapsto-N$, the mean curvature changes sign and the Gaussian curvature
is multiplied by $(-1)^n$.

\section{Computations in Semigeodesic Coordinates}

Let $(x^1,\ldots,x^n,v)$ be semigeodesic coordinates on
$U\subseteq\widetilde M$. For each $a$ with
$M_a=v^{-1}(a)\ne\varnothing$, the level set $M_a$ is a properly embedded
hypersurface. Its induced metric is denoted by $g_a$, and
\begin{equation}
\widetilde g=dv^2+sum_{\alpha,\beta=1}^n
g_{\alpha\beta}(x^1,\ldots,x^n,v)\,dx^\alpha dx^\beta. \tag{8.13}
\end{equation}
The restrictions of $(x^1,\ldots,x^n)$ are coordinates on $M_a$,
$g_a=g_{\alpha\beta}dx^\alpha dx^\beta|_{v=a}$, and $\partial_v$ is a unit
normal along $M_a$.

\begin{proposition}
\label{prop:semigeodesic-coordinates-second-form}
\dcref{ch8:8.17}
With respect to the normal $\partial_v$, the scalar second fundamental form,
shape operator, and mean curvature of $M_a$ have components
\[
(h_a)_{\alpha\beta}
=-\frac12\partial_vg_{\alpha\beta}|_{v=a},
\qquad
(s_a)_\beta{}^\alpha
=-\frac12g^{\alpha\gamma}\partial_vg_{\gamma\beta}|_{v=a},
\qquad
H_a=-\frac1{2n}g^{\alpha\beta}\partial_vg_{\alpha\beta}|_{v=a}.
\]
\end{proposition}

\begin{proof}
\sketch In semigeodesic coordinates, the normal component of
$\widetilde\nabla_{\partial_\alpha}\partial_\beta$ is
\[
\widetilde\Gamma^{n+1}_{\alpha\beta}\partial_v
=-\frac12\partial_vg_{\alpha\beta}\,\partial_v.
\]
Equation (8.9) gives the formula for $h_a$ after restriction to $v=a$.
Raising an index with $(g^{\alpha\gamma})$ gives $s_a$, and taking
$n^{-1}$ times its trace gives $H_a$.
\end{proof}

\section{Minimal Hypersurfaces}

For a compact embedded hypersurface $M$ with nonempty boundary, \emph{area}
means its induced $n$-dimensional volume. It is \emph{area-minimizing} if its
area is least among compact embedded hypersurfaces with the same boundary.

\begin{theorem}
\label{thm:area-minimizing-zero-mean-curvature}
\uses{ex:fermi-coordinates-hypersurface, thm:minimizing-geodesic-unit-speed}
\dcref{ch8:8.18}

Let $M$ be a compact codimension-one submanifold with nonempty boundary in a
Riemannian $(n+1)$-manifold $(\widetilde M,\widetilde g)$. If $M$ is
area-minimizing, then its mean curvature vanishes identically.
\end{theorem}

\begin{proof}
\sketch Fix $p\in\operatorname{Int}M$ and choose Fermi coordinates
$(x^1,\ldots,x^n,v)$ on $U\subseteq\widetilde M$ such that
$\widetilde U=U\cap M$ is a coordinate ball disjoint from $\partial M$.
For arbitrary $\varphi\in C_c^\infty(\widetilde U)$ and sufficiently small
$t$, replace $M\cap U$ by the graph $v=t\varphi(x)$ to obtain $M_t$. The graph
map
\[
f_t(x)=(x^1,\ldots,x^n,t\varphi(x))
\]
extends by the identity to a diffeomorphism $F_t\colon M\to M_t$. Pull the
induced metric back to $M$ and denote it by $g_t$. On $\widetilde U$,
\begin{equation}
(g_t)_{\alpha\beta}
=t^2\partial_\alpha\varphi\,\partial_\beta\varphi
+g_{\alpha\beta}(x,t\varphi(x)), \tag{8.15}
\end{equation}
and $g_t=g$ outside $\widetilde U$.

Since the area has a minimum at $t=0$, compute its first variation:
\begin{equation}
\left.\frac d{dt}\right|_{t=0}\operatorname{Area}(M_t)
=\int_{\widetilde U}\frac12(\det g)^{-1/2}
 \left.\frac\partial{\partial t}\right|_{t=0}\det g_t
 \,dx^1\cdots dx^n. \tag{8.16}
\end{equation}
Writing $\operatorname{cof}^{\alpha\beta}$ for the matrix cofactors gives
\begin{equation}
\left.\frac\partial{\partial t}\right|_{t=0}\det g_t
=\operatorname{cof}^{\alpha\beta}
 \left.\frac\partial{\partial t}\right|_{t=0}(g_t)_{\alpha\beta}. \tag{8.17}
\end{equation}
By Cramer's rule and (8.15), this equals
\[
(\det g)g^{\alpha\beta}\partial_vg_{\alpha\beta}\,\varphi.
\]
The semigeodesic formula for $H$ therefore yields
\begin{equation}
\left.\frac d{dt}\right|_{t=0}\operatorname{Area}(M_t)
=-n\int_{\widetilde U}H\varphi\,dV_g. \tag{8.18}
\end{equation}
The left side is zero for every compactly supported $\varphi$. A nonnegative
bump function supported where $H$ has either fixed positive or fixed negative
sign rules out $H(p)\ne0$. Thus $H=0$ on the interior, and continuity gives
$H=0$ on all of $M$.
\end{proof}

An immersed or embedded hypersurface, with or without boundary, is
\emph{minimal} if $H\equiv0$. This condition means stationarity for area, not
global area minimization, although sufficiently small pieces of a minimal
hypersurface are area-minimizing.

\begin{theorem}
\label{thm:constant-mean-curvature-isoperimetric}
\uses{thm:area-minimizing-zero-mean-curvature, thm:tubular-neighborhood}
\dcref{ch8:8.19}

Let $(\widetilde M,\widetilde g)$ be a Riemannian $(n+1)$-manifold, let
$D\subseteq\widetilde M$ be a compact regular domain, and set $M=\partial D$.
If $M$ minimizes surface area among boundaries of compact regular domains with
the same volume as $D$, then its mean curvature with respect to the outward
unit normal is constant.
\end{theorem}

\begin{proof}
\sketch Suppose $H$ is not constant and choose $p,q\in M$ with
$H(p)<H(q)$. By compactness and the tubular-neighborhood theorem, choose
disjoint ambient Fermi coordinate neighborhoods $\widetilde U$ of $p$ and
$\widetilde W$ of $q$, and set
$U=M\cap\widetilde U$ and $W=M\cap\widetilde W$. Choose normal coordinates
$v,w$ increasing outward, so locally $D$ is given by $v\leq0$ and $w\leq0$.
Shrink the neighborhoods and choose constants
\[
H(p)<H_1<H_2<H(q),
\qquad H\leq H_1\text{ on }U,
\qquad H\geq H_2\text{ on }W.
\]
Choose nonnegative functions
$\varphi\in C_c^\infty(U)$ and $\psi\in C_c^\infty(W)$ normalized by
\[
\int_U\varphi\,dV_g=\int_W\psi\,dV_g=1.
\]

For small $s,t$, perturb the two boundary patches to the graphs
$v=s\varphi$ and $w=t\psi$ by setting
\begin{align*}
D_{s,t}
&=\{z\in\widetilde U:v(z)\leq s\varphi(x(z))\}\\
&\quad\cup\{z\in\widetilde W:w(z)\leq t\psi(y(z))\}\\
&\quad\cup\bigl(D\setminus(\widetilde U\cup\widetilde W)\bigr),
\qquad M_{s,t}=\partial D_{s,t}.
\end{align*}
Put
\[
V(s,t)=\operatorname{Vol}(D_{s,t}),
\qquad A(s,t)=\operatorname{Area}(M_{s,t}).
\]
The first-variation calculation (8.18) gives
\[
\frac{\partial A}{\partial s}(0,0)
=-n\int_UH\varphi\,dV_g,
\qquad
\frac{\partial A}{\partial t}(0,0)
=-n\int_WH\psi\,dV_g.
\]
The fundamental theorem of calculus in the normal coordinates gives
\[
\frac{\partial V}{\partial s}(0,0)
=\int_U\varphi\,dV_g=1,
\qquad
\frac{\partial V}{\partial t}(0,0)
=\int_W\psi\,dV_g=1.
\]
The implicit function theorem supplies a smooth $\lambda(s)$ with
$V(s,\lambda(s))=V(0,0)$ and $\lambda'(0)=-1$. Area minimality under this
volume-preserving variation implies
\[
0=\left.\frac d{ds}\right|_{s=0}A(s,\lambda(s))
=-n\int_UH\varphi\,dV_g+n\int_WH\psi\,dV_g.
\]
But nonnegativity and normalization of the test functions give
\[
\int_UH\varphi\,dV_g
\leq H_1<H_2
\leq\int_WH\psi\,dV_g,
\]
a contradiction.
\end{proof}

For further treatment of minimal and constant-mean-curvature hypersurfaces,
see \cite{CM11}.
\section{Hypersurfaces in Euclidean Space}

Let $M\subseteq\mathbb R^{n+1}$ be an embedded hypersurface with induced
metric $g$, and let $h$ be its scalar second fundamental form for a chosen unit
normal.  Since the ambient Euclidean curvature vanishes, the Gauss and Codazzi
equations are
\begin{equation}
\tfrac12h\owedge h=\mathrm{Rm}, \tag{8.19}
\end{equation}
\begin{equation}
Dh=0, \tag{8.20}
\end{equation}
or, in a local frame,
\begin{equation}
h_{il}h_{jk}-h_{ik}h_{jl}=R_{ijkl}, \tag{8.21}
\end{equation}
\begin{equation}
h_{ij;k}-h_{ik;j}=0. \tag{8.22}
\end{equation}
A symmetric $2$-tensor satisfying $Dh=0$ is a \emph{Codazzi tensor}.

Conversely, (8.19)--(8.20) are locally sufficient: if $g$ is a Riemannian
metric and $h$ is symmetric and satisfies these equations, then every point
has a neighborhood admitting an isometric immersion into $\mathbb R^{n+1}$
whose scalar second fundamental form is $h$.

For $v\in T_pM$ with $|v|=1$, let $\gamma_v$ be the intrinsic geodesic with
initial velocity $v$.  Its Euclidean acceleration is
\[
\gamma_v''(0)=h(v,v)N_p.
\]
Thus $h(v,v)=\langle\gamma_v''(0),N_p\rangle$, with sign determined by the
chosen normal.

\begin{proposition}[Osculating Circle]
\label{prop:osculating-circle}
\dcref{ch8:8.21}
Let $\gamma:I\to\mathbb R^m$ be unit speed, let $t_0\in I$, and assume
$\kappa(t_0)\ne0$.
\begin{enumerate}
\item[(a)] There is a unique unit-speed circle
$c:\mathbb R\to\mathbb R^m$ whose position, velocity, and acceleration at
$t_0$ agree with those of $\gamma$.  It is the \emph{osculating circle} at
$\gamma(t_0)$.
\item[(b)] If its radius is $R$, then $\kappa(t_0)=1/R$.
\end{enumerate}
\end{proposition}

\begin{proof}
\sketch
For orthonormal $v,w\in\mathbb R^m$, the circle with center $q$ and radius
$R$ has unit-speed parametrization
\[
c(t)=q+R\cos\left(\frac{t-t_0}{R}\right)v
+R\sin\left(\frac{t-t_0}{R}\right)w,
\]
and
\[
c(t_0)=q+Rv,
\qquad c'(t_0)=w,
\qquad c''(t_0)=-\frac1Rv.
\]
Taking
\[
R=\frac1{|\gamma''(t_0)|}=\frac1{\kappa(t_0)},
\quad v=-R\gamma''(t_0),
\quad w=\gamma'(t_0),
\quad q=\gamma(t_0)-Rv
\]
gives the required circle and radius.
\end{proof}

\section{Computations in Euclidean Space}

Let $X:U\to M\subseteq\mathbb R^{n+1}$ be a local parametrization, with
$X_i=\partial_iX$.  The vectors $X_1,\ldots,X_n$ span $TM$, and Gram--Schmidt
against any independent ambient coordinate vector produces either choice of
unit normal $N$.

\begin{proposition}[Coordinate Formula for the Second Fundamental Form]
\label{prop:second-fundamental-form-coordinate-formula}
\dcref{ch8:8.23}
For a unit normal field $N$ along the parametrized hypersurface,
\begin{equation}
h(X_i,X_j)=\left\langle
\frac{\partial^2X}{\partial u^i\partial u^j},N\right\rangle. \tag{8.23}
\end{equation}
\end{proposition}

\begin{proof}
\sketch
The identity $\langle X_j,N\rangle=0$ implies
\[
0=\left\langle\partial_i\partial_jX,N\right\rangle
+\left\langle X_j,\overline\nabla_{X_i}N\right\rangle.
\]
The Weingarten equation gives
$\overline\nabla_{X_i}N=-s(X_i)$, so the second term is
$-h(X_j,X_i)$, proving (8.23).
\end{proof}

\begin{proposition}[Principal Curvatures and Local Shape]
\label{prop:principal-curvatures-local-shape}
\dcref{ch8:8.24}
Let $M\subset\mathbb R^{n+1}$ be a hypersurface, let $p\in M$, and let
$k_1,\ldots,k_n$ be its principal curvatures at $p$ for a chosen unit normal.
Some Euclidean isometry sends $p$ to $0$ and a neighborhood of $p$ in $M$ to
a graph $x^{n+1}=f(x^1,\ldots,x^n)$ with
\begin{equation}
f(x)=\frac12\left(k_1(x^1)^2+\cdots+k_n(x^n)^2\right)+O(|x|^3). \tag{8.24}
\end{equation}
\end{proposition}

\begin{proof}
\sketch
Translate and rotate so that $p=0$, $T_pM$ is the first-coordinate
hyperplane, the selected normal is $\partial_{n+1}$, and the second
fundamental form is diagonal with entries $k_1,\ldots,k_n$.  The implicit
function theorem represents $M$ locally by
$X(u)=(u^1,\ldots,u^n,f(u))$, with $f(0)=0$ and $df(0)=0$.  Formula (8.23)
then gives
\[
\frac{\partial^2f}{\partial x^i\partial x^j}(0)=h_{ij}(0).
\]
Taylor's theorem yields (8.24).
\end{proof}

If a smooth ambient field $N=N^i\partial_i$ restricts to a unit normal along
$M$, the Weingarten equation gives, for
$X=X^j\partial_j\in TM$,
\begin{equation}
sX=-\overline\nabla_XN
=-\sum_{i,j=1}^{n+1}X^j(\partial_jN^i)\partial_i. \tag{8.25}
\end{equation}
If $M$ is locally a regular level set of $F$, one may take
\[
N=\frac{\operatorname{grad}F}{|\operatorname{grad}F|}.
\]

\begin{example}[Shape Operators of Spheres]
\label{ex:shape-operators-spheres}
\uses{prop:existence-local-defining-functions}
\dcref{ch8:8.25}

For $F(x)=|x|^2$, the outward unit normal on $S^n(R)$ is
\[
N=\frac1R\sum_{i=1}^{n+1}x^i\partial_i.
\]
Therefore
\[
sX=-\frac1R\sum_{i,j=1}^{n+1}X^j(\partial_jx^i)\partial_i
=-\frac1RX.
\]
Thus every principal curvature and the mean curvature equal $-1/R$, while
the Gaussian curvature is $(-1/R)^n$ for this orientation.
\end{example}

For a parametrized surface in $\mathbb R^3$, a unit normal is
\begin{equation}
N=\frac{X_1\times X_2}{|X_1\times X_2|}. \tag{8.26}
\end{equation}

\section{The Gaussian Curvature of a Surface Is Intrinsic}

\begin{theorem}[Gauss's Theorema Egregium]
\label{thm:theorema-egregium}
\uses{cor:curvature-tensor-dim-2}
\dcref{ch8:8.27}

For an embedded Riemannian surface $(M,g)\subset\mathbb R^3$, the Gaussian
curvature satisfies
\[
K=\frac12S.
\]
Consequently $K$ is an intrinsic local-isometry invariant.
\end{theorem}

\begin{proof}
\sketch
At $p$, choose an orthonormal basis $(b_1,b_2)$ of $T_pM$.  The shape
operator and $h$ have the same matrix, and the Gauss equation gives
\[
\mathrm{Rm}_p(b_1,b_2,b_2,b_1)
=h_{11}h_{22}-h_{12}h_{21}=\det(h_{ij})=K(p).
\]
In dimension two,
$\mathrm{Rm}=\tfrac14Sg\owedge g$, so evaluation on
$(b_1,b_2,b_2,b_1)$ gives $\mathrm{Rm}_p(b_1,b_2,b_2,b_1)=\tfrac12S(p)$.
\end{proof}

For an abstract Riemannian $2$-manifold, define its \emph{Gaussian curvature}
by $K=\tfrac12S$.  This agrees with the determinant of the shape operator for
embedded surfaces in $\mathbb R^3$.

\begin{corollary}[Curvature Tensors of a Riemannian Surface]
\label{cor:riemann-two-manifold-curvatures}
\uses{cor:curvature-tensor-dim-2}
\dcref{ch8:8.28}

For every Riemannian $2$-manifold,
\[
\mathrm{Rm}=\frac12K g\owedge g,
\qquad
\mathrm{Rc}=Kg,
\qquad
S=2K.
\]
\end{corollary}

\section{Sectional Curvatures}

Let $(M,g)$ be a Riemannian $n$-manifold with $n\geq2$, let $p\in M$, and
let $V\subseteq T_pM$ be a star-shaped neighborhood on which $\exp_p$ is a
diffeomorphism onto $U$.  For a $2$-plane $\Pi\subseteq T_pM$, the
\emph{plane section}
\[
S_\Pi=\exp_p(\Pi\cap V)
\]
is an embedded surface through $p$ with $T_pS_\Pi=\Pi$.  Its intrinsic
Gaussian curvature at $p$ is the \emph{sectional curvature} $\sec(\Pi)$.
For a basis $(v,w)$ of $\Pi$, write $\sec(v,w)=\sec(\Pi)$.

For vectors in an inner-product space, set
\begin{equation}
|v\wedge w|=\sqrt{|v|^2|w|^2-\langle v,w\rangle^2}. \tag{8.27}
\end{equation}
This quantity is positive exactly when $v,w$ are independent and equals $1$
for an orthonormal pair.

\begin{proposition}[Formula for the Sectional Curvature]
\label{prop:sectional-curvature-formula}
\dcref{ch8:8.29}
If $v,w\in T_pM$ are linearly independent, then
\[
\sec(v,w)=\frac{\mathrm{Rm}_p(v,w,w,v)}{|v\wedge w|^2}.
\]
\end{proposition}

\begin{proof}
\sketch
Let $\Pi=\operatorname{span}(v,w)$.  If $z\in\Pi$, the ambient geodesic
$\gamma_z(t)=\exp_p(tz)$ lies in $S_\Pi$ near $0$.  The Gauss formula along
$\gamma_z$ gives
\[
0=D_t\gamma_z'=\overline D_t\gamma_z'+\mathrm{II}(\gamma_z',\gamma_z'),
\]
whose tangent and normal terms vanish separately.  Hence
$\mathrm{II}_p(z,z)=0$ for all $z\in\Pi$, and polarization gives
$\mathrm{II}_p=0$.  The Gauss equation therefore identifies the ambient and
surface curvature tensors on $\Pi$ at $p$.  For an orthonormal basis
$(b_1,b_2)$ of $\Pi$,
\[
\sec(\Pi)=\mathrm{Rm}_p(b_1,b_2,b_2,b_1).
\]

Apply Gram--Schmidt to an arbitrary basis:
\[
b_1=\frac v{|v|},
\qquad
b_2=\frac{w-\langle w,b_1\rangle b_1}
{|w-\langle w,b_1\rangle b_1|}.
\]
The curvature symmetries give
\begin{equation}
\sec(\Pi)=
\frac{\mathrm{Rm}_p(v,w,w,v)}
{|v|^2|w-\langle w,b_1\rangle b_1|^2}. \tag{8.29}
\end{equation}
The denominator is
$|v|^2|w|^2-\langle v,w\rangle^2=|v\wedge w|^2$.
\end{proof}

\begin{proposition}[Sectional Curvature Determines the Curvature Tensor]
\label{prop:sectional-curvature-determines-curvature-tensor}
\dcref{ch8:8.31}
Let $R_1,R_2$ be algebraic curvature tensors on a finite-dimensional
inner-product space $V$.  If
\[
\frac{R_1(v,w,w,v)}{|v\wedge w|^2}
=\frac{R_2(v,w,w,v)}{|v\wedge w|^2}
\]
for every linearly independent $v,w$, then $R_1=R_2$.
\end{proposition}

\begin{proof}
\sketch
Set $D=R_1-R_2$.  Then $D(v,w,w,v)=0$ for all $v,w$, by the hypothesis for
independent pairs and by curvature symmetries for dependent pairs.  Expanding
$D(v+w,x,x,v+w)=0$ gives
\[
D(v,x,x,w)=0.
\]
Expanding $D(v,x+u,x+u,w)=0$ then gives
\[
D(v,x,u,w)+D(v,u,x,w)=0,
\]
so $D$ is antisymmetric in every adjacent pair.  The algebraic Bianchi
identity now yields
\[
0=D(v,w,x,u)+D(w,x,v,u)+D(x,v,w,u)=3D(v,w,x,u),
\]
hence $D=0$.
\end{proof}

\begin{proposition}[Geometric Interpretation of Ricci and Scalar Curvatures]
\label{prop:geometric-interpretation-ricci-scalar-curvatures}
\uses{lem:symmetric-bilinear-form-unit-vectors}
\dcref{ch8:8.32}

Let $(M,g)$ be a Riemannian $n$-manifold and $p\in M$.
\begin{enumerate}
\item[(a)] If $v$ is unit and $(b_1,\ldots,b_n)$ is orthonormal with $b_1=v$,
then
\[
\mathrm{Rc}_p(v,v)=\sum_{k=2}^n\sec(b_1,b_k).
\]
\item[(b)] For any orthonormal basis,
\[
S(p)=\sum_{j\ne k}\sec(b_j,b_k),
\]
the sum over ordered pairs of distinct basis vectors.
\end{enumerate}
\end{proposition}

\begin{proof}
\sketch
Tracing the curvature tensor in the adapted orthonormal basis gives
\[
\mathrm{Rc}_p(v,v)
=\sum_{k=1}^n\mathrm{Rm}_p(b_k,b_1,b_1,b_k)
=\sum_{k=2}^n\sec(b_1,b_k).
\]
Tracing once more gives
\[
S(p)=\sum_{j=1}^n\mathrm{Rc}_p(b_j,b_j)
=\sum_{j\ne k}\sec(b_j,b_k).
\]
\end{proof}

Thus a common strict or weak sign for all sectional curvatures gives the same
sign for Ricci and scalar curvature.  The convention used here is
$\sec(v,w)=\mathrm{Rm}(v,w,w,v)/|v\wedge w|^2$; reversing the curvature-tensor
sign convention requires the corresponding reversed sign/order in this
formula.

\section{Sectional Curvatures of the Model Spaces}

A Riemannian manifold has \emph{constant sectional curvature} if all tangent
$2$-planes at all points have the same sectional curvature.

\begin{lemma}[Frame Homogeneity Implies Constant Sectional Curvature]
\label{lem:frame-homogeneous-constant-sectional-curvature}
\dcref{ch8:8.33}
Every frame-homogeneous Riemannian manifold has constant sectional curvature.
\end{lemma}

\begin{proof}
\sketch
Frame homogeneity supplies an isometry carrying any chosen tangent $2$-plane
to any other; sectional curvature is invariant under isometries.
\end{proof}

\begin{theorem}[Sectional Curvatures of the Model Spaces]
\label{thm:sectional-curvatures-model-spaces}
\uses{ex:shape-operators-spheres}
\dcref{ch8:8.34}

The model spaces have constant sectional curvatures
\begin{enumerate}
\item[(a)] $0$ on $(\mathbb R^n,\bar g)$;
\item[(b)] $1/R^2$ on $(\mathbb S^n(R),\bar g_R)$;
\item[(c)] $-1/R^2$ on $(\mathbb H^n(R),\bar g_R)$.
\end{enumerate}
\end{theorem}

\begin{proof}
\sketch
Euclidean space has zero curvature tensor, hence zero sectional curvature.
For the round sphere, the shape operator is $-(1/R)\operatorname{Id}$, so the
Gauss equation gives sectional curvature $(-1/R)^2=1/R^2$; frame homogeneity
makes this value constant.
\end{proof}

For each $c\in\mathbb R$, exactly one of these model families has constant
sectional curvature $c$, with its radius determined by $|c|^{-1/2}$ when
$c\ne0$.

\begin{proposition}[Metrics of Constant Sectional Curvature]
\label{prop:constant-sectional-curvature}
\dcref{ch8:8.36}
A Riemannian metric $g$ has constant sectional curvature $c$ if and only if
\[
\mathrm{Rm}=\frac12c\,g\owedge g.
\]
In this case,
\[
\mathrm{Rc}=(n-1)cg,
\qquad
S=n(n-1)c,
\]
\[
R(v,w)x=c\bigl(\langle w,x\rangle v-\langle v,x\rangle w\bigr),
\]
and, in any basis,
\[
R_{ijkl}=c(g_{il}g_{jk}-g_{ik}g_{jl}),
\qquad
R_{ij}=(n-1)cg_{ij}.
\]
\end{proposition}
\section{Further Results}

\begin{proposition}
\label{prop:complex-projective-space-curvature}
\dcref{ch8:8-13}
\uses{prop:oneill-formula}

Let $p:S^{2n+1}\to\mathbb{CP}^n$ be the Riemannian submersion of
Example~2.30. Identify $\mathbb C^{n+1}$ with $\mathbb R^{2n+2}$ using
$z^j=x^j+iy^j$.
\begin{enumerate}
\item[(a)] The field
\[
S=x^j\frac{\partial}{\partial y^j}
-y^j\frac{\partial}{\partial x^j}
\]
is tangent to $S^{2n+1}$ and spans the vertical space $V_z$ at each
$z\in S^{2n+1}$, with $j$ summed from $1$ to $n+1$.
\item[(b)] For horizontal vector fields $W,Z$ on $S^{2n+1}$,
\[
\begin{gathered}
[W,Z]^V=-d\omega(W,Z)S\\
=2\langle W,JZ\rangle S,
\end{gathered}
\]
where
\[
\omega=S^\flat=\sum_j(x^j\,dy^j-y^j\,dx^j)
\]
and the real-linear orthogonal map $J$ is
\[
\begin{gathered}
J\left(a^j\frac{\partial}{\partial x^j}
+b^j\frac{\partial}{\partial y^j}\right)\\
=a^j\frac{\partial}{\partial y^j}
-b^j\frac{\partial}{\partial x^j}.
\end{gathered}
\]
Thus $J$ is complex multiplication by $i$ and $J^2=-\operatorname{Id}$.
\item[(c)] O'Neill's formula gives
\begin{align*}
&Rm(w,x,y,z)\\
&=\langle\widetilde w,\widetilde z\rangle
  \langle\widetilde x,\widetilde y\rangle
\\
&\quad-\langle\widetilde w,\widetilde y\rangle
  \langle\widetilde x,\widetilde z\rangle\\
&\quad-2\langle\widetilde w,J\widetilde x\rangle
  \langle\widetilde y,J\widetilde z\rangle
\\
&\quad-\langle\widetilde w,J\widetilde y\rangle
  \langle\widetilde x,J\widetilde z\rangle\\
&\quad+\langle\widetilde w,J\widetilde z\rangle
  \langle\widetilde x,J\widetilde y\rangle
\end{align*}
for $q\in\mathbb{CP}^n$ and $w,x,y,z\in T_q\mathbb{CP}^n$, where the tilded
vectors are horizontal lifts to any
$\widetilde q\in p^{-1}(q)\subseteq S^{2n+1}$.
\item[(d)] For orthonormal $w,x\in T_q\mathbb{CP}^n$,
\[
\sec(w,x)=1+3\langle\widetilde w,J\widetilde x\rangle^2.
\]
\end{enumerate}
\end{proposition}

\begin{exercise}
\label{prob:complex-projective-space-curvature}
\dcref{ch8:8-13}
Let $p: S^{2n+1} \to \mathbb{CP}^n$ be the Riemannian submersion described in Example 2.30. In this problem, we identify $\mathbb{C}^{n+1}$ with $\mathbb{R}^{2n+2}$ by means of coordinates $(x^1, y^1, \ldots, x^{n+1}, y^{n+1})$ defined by $z^j = x^j + iy^j$.

(a) Show that the vector field
\[
S = x^j \frac{\partial}{\partial y^j} - y^j \frac{\partial}{\partial x^j}
\]
on $\mathbb{C}^{n+1}$ is tangent to $S^{2n+1}$ and spans the vertical space $V_z$ at each point $z \in S^{2n+1}$. (The implicit summation here is from 1 to $n+1$.)

(b) Show that for all horizontal vector fields $W, Z$ on $S^{2n+1}$,
\[
[W, Z]^V = -d\omega(W, Z)S = 2\langle W, JZ \rangle S,
\]
where $\omega$ is the 1-form on $\mathbb{C}^{n+1}$ given by
\[
\omega = S^\flat = \sum_j x^j dy^j - y^j dx^j,
\]
and $J: T\mathbb{C}^{n+1} \to T\mathbb{C}^{n+1}$ is the real-linear orthogonal map given by
\[
J\left(a^j \frac{\partial}{\partial x^j} + b^j \frac{\partial}{\partial y^j}\right) = a^j \frac{\partial}{\partial y^j} - b^j \frac{\partial}{\partial x^j}.
\]
(This is just multiplication by $i = \sqrt{-1}$ in complex coordinates. Notice that $J \circ J = -\text{Id}$.)

(c) Using O'Neill's formula (Problem 7-14), show that the curvature tensor of $\mathbb{CP}^n$ satisfies
\[
\text{Rm}(w, x, y, z) = \langle \tilde{w}, \tilde{z} \rangle \langle \tilde{x}, \tilde{y} \rangle - \langle \tilde{w}, \tilde{y} \rangle \langle \tilde{x}, \tilde{z} \rangle
- 2\langle \tilde{w}, J\tilde{x} \rangle \langle \tilde{y}, J\tilde{z} \rangle - \langle \tilde{w}, J\tilde{y} \rangle \langle \tilde{x}, J\tilde{z} \rangle
+ \langle \tilde{w}, J\tilde{z} \rangle \langle \tilde{x}, J\tilde{y} \rangle,
\]
for every $q \in \mathbb{CP}^n$ and $w, x, y, z \in T_q\mathbb{CP}^n$, where $\tilde{w}, \tilde{x}, \tilde{y}, \tilde{z}$ are horizontal lifts of $w, x, y, z$ to an arbitrary point $\tilde{q} \in p^{-1}(q) \subseteq S^{2n+1}$.

(d) Using the notation of part (c), show that for orthonormal vectors $w, x \in T_q\mathbb{CP}^n$, the sectional curvature of the plane spanned by $\{w, x\}$ is
\[
\sec(w, x) = 1 + 3\langle \tilde{w}, J\tilde{x} \rangle^2.
\]
\end{exercise}


\begin{proposition}
\label{prop:curvature-operator}
\dcref{ch8:8-33}
\uses{prop:complex-projective-space-curvature}

Let $(M,g)$ be Riemannian, and let $\Lambda^2(TM)$ be its bundle of
alternating contravariant $2$-tensors.
\begin{enumerate}
\item[(a)] There is a unique fiber metric on $\Lambda^2(TM)$ whose
norm satisfies
\[
|w\wedge x|^2=|w|^2|x|^2-\langle w,x\rangle^2.
\]
\item[(b)] There is a unique bundle endomorphism
$\mathcal R:\Lambda^2(TM)\to\Lambda^2(TM)$ satisfying
\begin{equation}
\begin{aligned}
&\langle\mathcal R(w\wedge x),y\wedge z\rangle\\
&=-Rm(w,x,y,z)
\end{aligned}
\tag{8.30}
\end{equation}
for tangent vectors at a common point. This is the \emph{curvature operator}.
\item[(c)] Positive and negative definiteness of $\mathcal R$ imply positive
and negative sectional curvature, respectively.

\emph{Sign convention.} The minus sign in (8.30) is required by the curvature
convention used here; it is omitted when the Riemann tensor is defined with
the opposite sign.
\item[(d)] The converse need not hold: the Fubini--Study metric on
$\mathbb{CP}^2$ has positive sectional curvature, but its curvature operator
is not positive definite.
\end{enumerate}
\end{proposition}

\begin{exercise}
\label{prob:curvature-operator}
\dcref{ch8:8-33}
\uses{prob:complex-projective-space-curvature}
Suppose $(M, g)$ is a Riemannian manifold. Let $\Lambda^2(TM)$ denote the bundle of 2-vectors (alternating contravariant 2-tensors) on $M$.

(a) Show that there is a unique fiber metric on $\Lambda^2(TM)$ whose associated norm satisfies
\[
|w \wedge x|^2 = |w|^2|x|^2 - \langle w,x\rangle^2
\]
for all tangent vectors $w, x$ at every point $q \in M$.

(b) Show that there is a unique bundle endomorphism $\mathcal{R}: \Lambda^2(TM) \to \Lambda^2(TM)$, called the \textbf{curvature operator} of $g$, that satisfies
\begin{equation}
\langle \mathcal{R}(w \wedge x), y \wedge z\rangle = -\text{Rm}(w,x,y,z)
\tag{8.30}
\end{equation}
for all tangent vectors $w, x, y, z$ at a point of $M$, where the inner product on the left-hand side is the one described in part (a).

(c) A Riemannian metric is said to have \textbf{positive curvature operator} if $\mathcal{R}$ is positive definite, and \textbf{negative curvature operator} if $\mathcal{R}$ is negative definite. Show that positive curvature operator implies positive sectional curvature, and negative curvature operator implies negative sectional curvature. [Remark: This is the reason for the negative sign in (8.30). If the Riemann curvature tensor is defined as the negative of ours, the negative sign should be omitted.]

(d) Show that the converse need not be true, by using the results of Problem 8-13 to compute the following expression on $\mathbb{CP}^2$ with the Fubini--Study metric:
\[
\langle \mathcal{R}(w \wedge x + y \wedge z), w \wedge x + y \wedge z\rangle,
\]
where $w, x, y, z$ are orthonormal vectors at an arbitrary point of $\mathbb{CP}^2$, chosen so that their horizontal lifts satisfy $\bar{y} = J\bar{w}$ and $\bar{z} = J\bar{x}$.
\end{exercise}


\begin{exercise}
\label{exer:codazzi-tensor-characterization}
\dcref{ch8:8.20}
Show that a smooth 2-tensor field $h$ on a Riemannian manifold is a Codazzi tensor if and only if both $h$ and $\nabla h$ are symmetric.
\end{exercise}

\begin{exercise}
\label{exer:connection-in-bundle}
\dcref{ch8:8.8}
Prove that $\nabla^F$ is a connection in $F$.
\end{exercise}

\begin{exercise}
\label{exer:constant-sectional-curvature-rp}
\dcref{ch8:8.35}
\uses{ex:real-projective-space-metric}
Show that the metric on real projective space $\mathbb{RP}^n$ defined in Example 2.34 has constant positive sectional curvature.
\end{exercise}

\begin{exercise}
\label{exer:lagrange-multiplier-eigenvector}
\dcref{ch8:8.15}
\uses{prop:lagrange-multiplier-rule}
Use the Lagrange multiplier rule (Prop. A.29) to prove this lemma.
\end{exercise}

\begin{exercise}
\label{exer:metric-scaling-sectional-curvature}
\dcref{ch8:8.30}
\uses{thm:conformal-transformation-curvature}
Suppose $(M,g)$ is a Riemannian manifold and $\tilde{g} = \lambda g$ for some positive constant $\lambda$. Use Theorem 7.30 to prove that for every $p \in M$ and plane $\Pi \subseteq T_p M$, the sectional curvatures of $\Pi$ with respect to $\tilde{g}$ and $g$ are related by $\sec_{\tilde{g}}(\Pi) = \lambda^{-1} \sec_g(\Pi)$.
\end{exercise}

\begin{exercise}
\label{exer:osculating-circle-uniqueness}
\dcref{ch8:8.22}
Complete the proof of the preceding proposition by proving uniqueness of the osculating circle.
\end{exercise}

\begin{exercise}
\label{exer:plane-cylinder-isometric}
\dcref{ch8:8.26}
Let $M_1 \subset \mathbb{R}^3$ be the plane $\{z = 0\}$, and let $M_2 \subset \mathbb{R}^3$ be the cylinder $\{x^2 + y^2 = 1\}$. Show that $M_1$ and $M_2$ are locally isometric, but the former has mean curvature zero, while the latter has mean curvature $\pm 1/2$, depending on which normal is chosen.
\end{exercise}

\begin{exercise}
\label{exer:prove-normal-connection}
\dcref{ch8:8.7}
\uses{prop:normal-connection-well-defined}
Prove the preceding proposition.
\end{exercise}

\begin{exercise}[The associated quadratic form]
\label{prob:associated-quadratic-form}
\dcref{ch8:8-21}

For every linear endomorphism $A : \mathbb{R}^n \to \mathbb{R}^n$, the \textbf{associated quadratic form} is the function $Q : \mathbb{R}^n \to \mathbb{R}$ defined by $Q(x) = \langle Ax, x \rangle$. Prove that
\[
\int_{S^{n-1}} Q \, dV_g = \frac{1}{n}(\operatorname{tr} A)\operatorname{Vol}(S^{n-1}).
\]

[Hint: Show that $\int_{S^{n-1}} x^i x^j \, dV_g = 0$ when $i \ne j$ by examining the effect of the isometry
\[
(x^1, \ldots, x^i, \ldots, x^j, \ldots, x^n) \mapsto (x^1, \ldots, x^j, \ldots, -x^i, \ldots, x^n),
\]
and compute $\int_{S^{n-1}} (x^i)^2 \, dV_g$ using the fact that $\sum_i (x^i)^2 = 1$ on the sphere.] [Remark: It is a standard fact of linear algebra that the trace of $A$ is independent of the choice of basis. This gives a geometric interpretation to the trace as $n$ times the average value of the associated quadratic form $Q$.]
(Used on p.\ 313.)
\end{exercise}

\begin{exercise}
\label{prob:berger-metric-su2}
\dcref{ch8:8-16}
For each $a > 0$, let $g_a$ be the Berger metric on $\text{SU}(2)$ defined in Problem 3-10. Compute the sectional curvatures with respect to $g_a$ of the planes spanned by $(X, Y)$, $(Y, Z)$, and $(Z, X)$. Prove that if $a \neq 1$, then $(\text{SU}(2), g_a)$ is homogeneous but not isotropic anywhere.
\end{exercise}

\begin{exercise}
\label{prob:catenoid-boundary-minimal}
\dcref{ch8:8-8}
\uses{prob:catenoid-minimal-surface}
For $h > 0$, let $C_h \subseteq \mathbb{R}^3$ be the one-dimensional submanifold $\{(x, y, z) : x^2 + y^2 = 1, z = \pm h\}$; it is the union of two unit circles lying in the $z = \pm h$ planes. Let $h_0 = 1/\sinh c$, where $c$ is the unique positive solution to the equation $c \tanh c = 1$. Prove that if $h < h_0$, then there are two distinct positive values $w_1 < w_2$ such that the portions of the catenoids $M_{w_1}$ and $M_{w_2}$ lying in the region $|z| \leq h$ are minimal surfaces with boundary $C_h$ (using the notation of Problem 8-7), while if $h = h_0$, there is exactly one such value, and if $h > h_0$, there are none. [Remark: This phenomenon can be observed experimentally. If you dip two parallel wire circles into soapy water, as long as they are close together they will form an area-minimizing soap film in the shape of a portion of a catenoid (the one with the larger ``waist,'' $M_{w_2}$ in the notation above). As you pull the circles apart, the ``waist'' of the catenoid gets smaller, until the ratio of the distance between the circles to their diameter reaches $h_0 \approx 0.6627$, at which point the film will burst and the only soap film that can be formed is two disjoint flat disks.]
\end{exercise}

\begin{exercise}
\label{prob:catenoid-minimal-surface}
\dcref{ch8:8-7}
For $w > 0$, let $M_w \subseteq \mathbb{R}^3$ be the surface of revolution obtained by revolving the curve $\gamma(t) = (w \cosh(t/w), t)$ around the $z$-axis, called a \textbf{catenoid}. Show that $M_w$ is a minimal surface for each $w$.
\end{exercise}

\begin{exercise}
\label{prob:convex-hypersurface}
\dcref{ch8:8-18}
Suppose $(M, g)$ is a Riemannian hypersurface in a Riemannian manifold $(\overline{M}, \overline{g})$, and $N$ is a unit normal vector field along $M$. We say that $M$ is \textbf{convex} (with respect to $N$) if its scalar second fundamental form satisfies $h(v, v) \leq 0$ for all $v \in TM$. Show that if $M$ is convex and $\overline{M}$ has sectional curvatures bounded below by a constant $c$, then all sectional curvatures of $M$ are bounded below by $c$.
\end{exercise}

\begin{exercise}
\label{prob:curvature-tensors-constant-curvature}
\dcref{ch8:8-29}
\uses{prop:constant-sectional-curvature}
Prove Proposition 8.36 (curvature tensors on constant-curvature spaces).
\end{exercise}

\begin{exercise}
\label{prob:einstein-sectional-curvature}
\dcref{ch8:8-14}
Show that a Riemannian 3-manifold is Einstein if and only if it has constant sectional curvature.
\end{exercise}

\begin{exercise}
\label{prob:fubini-study-not-conformal-flat}
\dcref{ch8:8-27}
\uses{prob:complex-projective-space-curvature,prob:weyl-tensor-4d}
Prove that the Fubini-Study metric on $\mathbb{CP}^2$ is not locally conformally flat.

[Hint: Use Problems 8-13 and 8-26.]
\end{exercise}

\begin{exercise}
\label{prob:gauss-map}
\dcref{ch8:8-9}
Let $M \subseteq \mathbb{R}^{n+1}$ be a Riemannian hypersurface, and let $N$ be a smooth unit normal vector field along $M$. At each point $p \in M$, $N_p \in T_p \mathbb{R}^{n+1}$ can be thought of as a unit vector in $\mathbb{R}^{n+1}$ and therefore as a point in $S^n$. Thus each choice of unit normal vector field defines a smooth map $\nu : M \to S^n$, called the \textbf{Gauss map} of $M$. Show that $\nu^* dV_g = (-1)^n K \, dV_g$, where $K$ is the Gaussian curvature of $M$.
\end{exercise}

\begin{exercise}
\label{prob:geodesic-curvature-non-unit-speed-curve}
\dcref{ch8:8-6}

Suppose $(M, g)$ is a Riemannian manifold, and $\gamma: I \to M$ is a regular (but not necessarily unit-speed) curve in $M$. Show that the geodesic curvature of $\gamma$ at $t \in I$ is
\[
\kappa(t) = \frac{|\gamma'(t) \times D_t \gamma'(t)|}{|\gamma'(t)|^3},
\]
where the norm in the numerator is the one defined by (8.27). Show also that in $\mathbb{R}^3$ with the Euclidean metric, the formula can be written
\[
\kappa(t) = \frac{|\gamma'(t) \times \gamma''(t)|}{|\gamma'(t)|^3}.
\]
\end{exercise}

\begin{exercise}
\label{prob:graph-submanifold-shape}
\dcref{ch8:8-1}
\uses{ex:metrics-in-graph-coordinates}
Suppose $U$ is an open set in $\mathbb{R}^n$ and $f: U \to \mathbb{R}$ is a smooth function. Let $M = \{(x, f(x)) : x \in U\} \subseteq \mathbb{R}^{n+1}$ be the graph of $f$, endowed with the induced Riemannian metric and upward unit normal (see Example 2.19).

\begin{enumerate}
\item[(a)] Compute the components of the shape operator in graph coordinates, in terms of $f$ and its partial derivatives.
\item[(b)] Let $M \subseteq \mathbb{R}^{n+1}$ be the $n$-dimensional paraboloid defined as the graph of $f(x) = |x|^2$. Compute the principal curvatures of $M$.
\end{enumerate}

(Used on p. 361.)
\end{exercise}

\begin{exercise}
\label{prob:homogeneous-isotropic-3-manifold}
\dcref{ch8:8-15}
\uses{prob:complex-projective-space-curvature}
Suppose $(M, g)$ is a 3-dimensional Riemannian manifold that is homogeneous and isotropic. Show that $g$ has constant sectional curvature. Show that the analogous result in dimension 4 is not true. [Hint: See Problem 8-13.]
\end{exercise}

\begin{exercise}
\label{prob:hyperbolic-constant-curvature}
\dcref{ch8:8-28}
\uses{thm:sectional-curvatures-model-spaces,thm:gauss-equation,prob:lorentz-hypersurface-formulas}
Complete the proof of Theorem 8.34 by showing in two ways that the hyperbolic space of radius $R$ has constant sectional curvature equal to $-1/R^2$.

\begin{enumerate}
\item[(a)] In the hyperboloid model, compute the second fundamental form of $\mathbb{H}^n(R) \subset \mathbb{R}^{n,1}$ at the point $(0,\ldots,0,R)$, and use either the general form of the Gauss equation (Thm.\ 8.5) or the formulas of Problem 8-19.

\item[(b)] In the Poincar\'e ball model, use formula (7.44) for the conformal transformation of the curvature to compute the Riemann curvature tensor at the origin.
\end{enumerate}
\end{exercise}

\begin{exercise}
\label{prob:hypersurface-defining-function}
\dcref{ch8:8-2}
Let $(M, g)$ be an embedded Riemannian hypersurface in a Riemannian manifold $(\bar{M}, \bar{g})$, let $F$ be a local defining function for $M$, and let $N = \grad F / |\grad F|$.

(a) Show that the scalar second fundamental form of $M$ with respect to the unit normal $N$ is given by
\[
h(X,Y) = -\frac{\bar{\nabla}^2 F(X,Y)}{|\operatorname{grad} F|}
\]
for all $X,Y \in \mathfrak{X}(M)$.

(b) Show that the mean curvature of $M$ is given by
\[
H = -\frac{1}{n} \operatorname{div}_{\bar{g}} \left( \frac{\operatorname{grad} F}{|\operatorname{grad} F|} \right),
\]
where $n = \dim M$ and $\operatorname{div}_{\bar{g}}$ is the divergence operator of $\bar{g}$ (see Problem 5-14). [Hint: First prove the following linear-algebra lemma: If $V$ is a finite-dimensional inner product space, $w \in V$ is a unit vector, and $A: V \to V$ is a linear map that takes $w^{\perp}$ to $w^{\perp}$, then $\operatorname{tr}(A|_{w^{\perp}}) = \operatorname{tr} B$, where $B: V \to V$ is defined by $Bx = Ax - \langle x, w \rangle Aw$.]
\end{exercise}

\begin{exercise}
\label{prob:k-point-homogeneous-ch8}
\dcref{ch8:8-24}
\uses{prob:k-point-homogeneous}
Let $(M,g)$ be a connected Riemannian manifold. Recall the definition of $k$-point homogeneous from Problem 6-18. Prove the following:
\begin{enumerate}
\item[(a)] If $(M,g)$ is homogeneous, then it has constant scalar curvature.
\item[(b)] If $(M,g)$ is 2-point homogeneous, then it is Einstein.
\item[(c)] If $(M,g)$ is 3-point homogeneous, then it has constant sectional curvature.
\end{enumerate}
\end{exercise}

\begin{exercise}
\label{prob:k-points-hyperbolic-geodesic}
\dcref{ch8:8-31}
For $1 < k \leq n$, show that any $k$ points in the hyperbolic space $\mathbb{H}^n(R)$ lie in a totally geodesic $(k-1)$-dimensional submanifold, which is isometric to $\mathbb{H}^{k-1}(R)$.
\end{exercise}

\begin{exercise}
\label{prob:lie-group-bi-invariant-metric-ch8}
\dcref{ch8:8-17}
Let $G$ be a Lie group with a bi-invariant metric $g$ (see Problem 7-13).

(a) Suppose $X$ and $Y$ are orthonormal elements of $\text{Lie}(G)$. Show that $\sec(X_p, Y_p) = \frac{1}{4}|[X, Y]|^2$ for each $p \in G$, and conclude that the sectional curvatures of $(G, g)$ are all nonnegative.

(b) Show that every Lie subgroup of $G$ is totally geodesic in $G$.

(c) Now suppose $G$ is connected. Show that $G$ is flat if and only if it is abelian.
\end{exercise}

\begin{exercise}
\label{prob:lie-group-fixed-point-geodesic}
\dcref{ch8:8-32}
Suppose $(M, g)$ is a Riemannian manifold and $G$ is a Lie group acting smoothly and isometrically on $M$. Let $S \subseteq M$ be the fixed point set of $G$, that is, the set of points $p \in M$ such that $\varphi(p) = p$ for all $\varphi \in G$. Show that each connected component of $S$ is a smoothly embedded totally geodesic submanifold of $M$. (The reason for restricting to a single connected component is that different components may have different dimensions.)
\end{exercise}

\begin{exercise}
\label{prob:lie-group-isometric-action-dimension}
\dcref{ch8:8-23}
\uses{prob:lie-group-dimension-bound}
Suppose $(M,g)$ is a connected $n$-dimensional Riemannian manifold, and a Lie group $G$ acts isometrically and effectively on $M$. Problem 5-12 showed that $\dim G \le n(n+1)/2$. Prove that if equality holds, then $g$ has constant sectional curvature.
\end{exercise}

\begin{exercise}
\label{prob:lorentz-einstein-field-equation}
\dcref{ch8:8-20}
\uses{prob:lorentz-hypersurface-formulas}
Suppose $(\overline{M}, \overline{g})$ is an $(n + 1)$-dimensional Lorentz manifold, and assume that $\overline{g}$ satisfies the Einstein field equation with a cosmological constant:
\[
\overline{Rc} - \frac{1}{2}\overline{S}\bar{g} + \Lambda\bar{g} = T,
\]
where $\Lambda$ is a constant and $T$ is a smooth symmetric 2-tensor field (see p.\ 211). Let $(M, g)$ be a Riemannian hypersurface in $\bar{M}$, and let $h$ be its scalar second fundamental form as defined in Problem 8-19. Use the results of Problem 8-19 to show that $g$ and $h$ satisfy the following \textbf{Einstein constraint equations} on $M$:
\[
S - 2\Lambda - |h|_g^2 + (\operatorname{tr}_g h)^2 = 2\rho,
\]
\[
\operatorname{div} h - d(\operatorname{tr}_g h) = J,
\]
where $S$ is the scalar curvature of $g$, $\operatorname{div} h$ the divergence of $h$ with respect to $g$ as defined in Problem 5-16, $\rho$ is the function $\rho = T(N, N)$, and $J$ is the 1-form $J(X) = T(N, X)$. [Remark: When $J$ and $\rho$ are both zero, these equations, known as the \textbf{vacuum Einstein constraint equations}, are necessary conditions for $g$ and $h$ to be the metric and second fundamental form of a Riemannian hypersurface in an $(n+1)$-dimensional Lorentz manifold satisfying the vacuum Einstein field equation with cosmological constant $\Lambda$. It was proved in 1950 by Yvonne Choquet-Bruhat that these conditions are also sufficient; for a proof, see [CB09, pp.\ 166--168].]
\end{exercise}

\begin{exercise}
\label{prob:lorentz-hypersurface-formulas}
\dcref{ch8:8-19}
\uses{thm:fundamental-equations-hypersurface}
Suppose $(M, g)$ is a Riemannian hypersurface in an $(n + 1)$-dimensional Lorentz manifold $(\overline{M}, \overline{g})$, and $N$ is a smooth unit normal vector field along $M$. Define the \textbf{scalar second fundamental form} $h$ and the \textbf{shape operator} $s$ by requiring that $\text{II}(X, Y) = h(X, Y)N$ and $\langle sX, Y \rangle = h(X, Y)$ for all $X, Y \in \mathfrak{X}(M)$. Prove the following Lorentz analogues of the formulas of Theorem 8.13 (with notation as in that theorem):

(a) \textsc{Gauss Formula:} $\overline{\nabla}_X Y = \nabla_X Y + h(X, Y)N$.

(b) \textsc{Gauss Formula for a Curve:} $\overline{D}_t X = D_t X + h(\gamma', X)N$.

(c) \textsc{Weingarten Equation:} $\overline{\nabla}_X N = sX$.

(d) \textsc{Gauss Equation:} $\overline{Rm}(W, X, Y, Z) = (Rm + \frac{1}{2} h \owedge h)(W, X, Y, Z)$.

(e) \textsc{Codazzi Equation:} $\overline{Rm}(W, X, Y, N) = -(Dh)(Y, W, X)$.
\end{exercise}

\begin{exercise}
\label{prob:paraboloid-isotropic-one-point}
\dcref{ch8:8-5}

Let $S \subseteq \mathbb{R}^3$ be the paraboloid given by $z = x^2 + y^2$, with the induced metric. Prove that $S$ is isotropic at only one point. (Used on p. 56.)
\end{exercise}

\begin{exercise}
\label{prob:product-metric-submanifolds}
\dcref{ch8:8-10}
\uses{exer:prove-lemma-2-11,prob:product-metric-curvature}
Suppose $g = g_1 \oplus g_2$ is a product metric on $M_1 \times M_2$ as in (2.12). (See also Problem 7-6.)

\begin{enumerate}
\item[(a)] Show that for each point $p_i \in M_i$, the submanifolds $M_1 \times \{p_2\}$ and $\{p_1\} \times M_2$ are totally geodesic.

\item[(b)] Show that $\sec(\Pi) = 0$ for every 2-plane $\Pi \subseteq T_{(p_1,p_2)}(M_1 \times M_2)$ spanned by $v_1 \in T_{p_1} M_1$ and $v_2 \in T_{p_2} M_2$.

\item[(c)] Show that if $M_1$ and $M_2$ both have nonnegative sectional curvature, then $M_1 \times M_2$ does too; and if $M_1$ and $M_2$ both have nonpositive sectional curvature, then $M_1 \times M_2$ does too.

\item[(d)] Now suppose $(M_i, g_i)$ has constant sectional curvature $c_i$ for $i = 1,2$; let $n_i = \dim M_i$. Show that $(M, g)$ is Einstein if and only if $c_1(n_1 - 1) = c_2(n_2 - 1)$; and $(M, g)$ has constant sectional curvature if and only if $n_1 = n_2 = 1$ or $c_1 = c_2 = 0$.
\end{enumerate}

(Used on p. 366.)
\end{exercise}

\begin{exercise}[Ricci curvature without choice of basis]
\label{prob:ricci-curvature-without-basis}
\dcref{ch8:8-22}
\uses{prob:associated-quadratic-form}

Let $(M, g)$ be a Riemannian $n$-manifold and $p \in M$. Proposition 8.32 gave a geometric interpretation of the Ricci curvature at $p$ based on a choice of orthonormal basis. This problem describes an interpretation that does not refer to a basis. For each unit vector $v \in T_p M$, prove that
\[
\operatorname{Rc}_p(v, v) = \frac{n-1}{\operatorname{Vol}(S_v^{\perp})} \int_{w \in S_v^{\perp}} \sec(v, w) \, d\bar{V}_g,
\]
where $S_v^{\perp}$ denotes the set of unit vectors in $T_p M$ that are orthogonal to $v$ and $\bar{g}$ denotes the Riemannian metric on $S_v^{\perp}$ induced from the flat metric $g_p$ on $T_p M$. [Hint: Use Problem 8-21.]
\end{exercise}

\begin{exercise}
\label{prob:ricci-tensor-hypersurface}
\dcref{ch8:8-30}
Suppose $M \subseteq \mathbb{R}^{n+1}$ is a hypersurface with the induced Riemannian metric. Show that the Ricci tensor of $M$ satisfies
\[
Rc(v,w) = \langle nHsv - s^2v, w\rangle,
\]
where $H$ and $s$ are the mean curvature and shape operator of $M$, respectively, and $s^2v = s(sv)$.
\end{exercise}

\begin{exercise}
\label{prob:riemannian-immersion-bounded-sectional-curvatures}
\dcref{ch8:8-12}
Suppose $\pi: (\tilde{M}, \tilde{g}) \to (M, g)$ is a Riemannian immersion, and $(\tilde{M}, \tilde{g})$ has all sectional curvatures bounded below by a constant $c$. Use O'Neill's formula (Problem 7-14) to show that the sectional curvatures of $(M, g)$ are bounded below by the same constant.
\end{exercise}

\begin{exercise}
\label{prob:surface-affine-plane-curvature}
\dcref{ch8:8-11}
Suppose $M \subseteq \mathbb{R}^3$ is a smooth surface, $p \in M$, and $N_p$ is a normal vector to $M$ at $p$. Show that if $\Pi \subseteq \mathbb{R}^3$ is any affine plane containing $p$ and parallel to $N_p$, then there is a neighborhood $U$ of $p$ in $\mathbb{R}^3$ such that $M \cap \Pi \cap U$ is a smooth embedded curve in $\Pi$, and the principal curvatures of $M$ at $p$ are equal to the minimum and maximum signed Euclidean curvatures of these curves as $\Pi$ varies. [Remark: This justifies the informal recipe for computing principal curvatures given in Chapter 1.]
\end{exercise}

\begin{exercise}
\label{prob:surface-revolution-constant-gaussian-curvature}
\dcref{ch8:8-4}
\uses{cor:local-uniqueness-constant-curvature}

Show that there is a surface of revolution in $\mathbb{R}^3$ that has constant Gaussian curvature equal to 1 but does not have constant principal curvatures. [Remark: We will see in Chapter 10 that this surface is locally isometric to $S^2$ (see Cor. 10.15), so this gives an example of two nonflat hypersurfaces in $\mathbb{R}^3$ that are locally isometric but have different principal curvatures.]
\end{exercise}

\begin{exercise}
\label{prob:surface-revolution-shape-principal-curvatures}
\dcref{ch8:8-3}
\uses{ex:surfaces-of-revolution}

Let $C$ be an embedded smooth curve in the half-plane $H = \{(r,z) : r > 0\}$, and let $S_C \subseteq \mathbb{R}^3$ be the surface of revolution determined by $C$ as in Example 2.20. Let $\gamma(t) = (a(t), b(t))$ be a unit-speed local parametrization of $C$, and let $X$ be the local parametrization of $S_C$ given by (2.11).

(a) Compute the shape operator and principal curvatures of $S_C$ in terms of $a$ and $b$, and show that the principal directions at each point are tangent to the meridians and latitude circles.

(b) Show that the Gaussian curvature of $S_C$ at a point $X(t,\theta)$ is equal to $-a''(t)/a(t)$.
\end{exercise}

\begin{exercise}
\label{prob:weyl-tensor-4d}
\dcref{ch8:8-26}
Let $(M,g)$ be a 4-dimensional Riemannian manifold. Given $p \in M$ and an orthonormal basis $(b_i)$ for $T_p M$, prove that the Weyl tensor at $p$ satisfies
\[
W_{1221} = \frac{1}{3}(k_{12} + k_{34}) - \frac{1}{6}(k_{13} + k_{14} + k_{23} + k_{24}),
\]
where $k_{ij}$ is the sectional curvature of the plane spanned by $(b_i, b_j)$.
\end{exercise}

\begin{exercise}
\label{prob:weyl-tensor-product-metric}
\dcref{ch8:8-25}
\uses{prob:product-metric-curvature}
For $i = 1,2$, suppose $(M_i, g_i)$ is a Riemannian manifold of dimension $n_i \ge 2$ with constant sectional curvature $c_i$; and let $g = g_1 \oplus g_2$ be the product metric on $M = M_1 \times M_2$. Show that the Weyl tensor of $g$ is given by
\[
W = \frac{c_1 + c_2}{2(n-1)(n-2)} \left( n_2(n_2-1)h_1 \owedge h_1 - 2(n_1-1)(n_2-1)h_1 \owedge h_2 + n_1(n_1-1)h_2 \owedge h_2 \right),
\]
where $n = n_1 + n_2$ and $h_i = \pi_i^* g_i$ for $i = 1,2$. Conclude that $(M,g)$ is locally conformally flat if and only if $c_2 = -c_1$. (See also Problem 7-6.)

(Used on p.~67.)
\end{exercise}
